\begin{figure}[htbp]
\centering
\vspace*{\fill}
\setbox\figbox=\hbox{
\begin{tikzpicture}[scale=1.75]
\path[use as bounding box] (-7,-5) rectangle (7,5);

% Angular sectors as wedges
\definecolor{color0}{RGB}{255,0,0}
\definecolor{color1}{RGB}{255,255,0}
\definecolor{color2}{RGB}{0,255,0}
\definecolor{color3}{RGB}{2,192,230}
\definecolor{color4}{RGB}{2,67,230}
\definecolor{color5}{RGB}{188, 2, 230}
\definecolor{nearestgreen}{RGB}{0,180,0}

\begin{scope}
\clip (-4,-4) rectangle (4,4);
\foreach \t in {0,...,5} {
  \fill[color=color\t, fill opacity=0.2] (0,0) -- (\t*60:10) arc (\t*60:(\t+1)*60:10) -- cycle;
}
\end{scope}

% Six dashed radial rays to demonstrate order-6 rotational symmetry (extended to 5 for coverage)
\foreach \t in {0,...,5} {
  \pgfmathsetmacro\rad{ifthenelse(\t==0 || \t==3, 4, 4.6)}
  \draw[dashed] (0,0) -- (\t*60:\rad);
}

% Real and imaginary axes (thin black solid lines, extended)
\draw[black,semithick,<->,>={Latex[scale=2]}] (-4,0) -- (4,0); % Real axis extended to cover max x ≈3.5
\draw[black,semithick,<->,>={Latex[scale=2]}] (0,-4) -- (0,4); % Imaginary axis

\draw[dashed,thick,color=nearestgreen] (0,0) circle (1);  % r=1 circle
% \node[font=\Large,color=blue] at (10:2.0) {$\Lambda_{-,\sqrt{1}}^4$};
% \draw[blue, solid, thick, ->] (1.55, 0.3) -> (0.65, 0.15);
% \node[font=\Large,color=nearestgreen] at (30:1.35) {$\Lambda_{T,\sqrt{1}}^4$};
% \node[font=\Large,color=red] at (30:2.35) {$\Lambda_{+,\sqrt{1}}^4$};

% Lattice vertices and edges (symmetric truncation within radius R=4)
\foreach \m in {-8,...,8} {
  \foreach \n in {-8,...,8} {  % Balanced n range to match m
    \pgfmathsetmacro\x{\m + \n * 0.5}
    \pgfmathsetmacro\y{\n * (sqrt(3)/2)}
    \pgfmathsetmacro\norm{sqrt(\x*\x + \y*\y)}
    \pgfmathparse{int(\m*\m + \m*\n + \n*\n)}  % Integer norm squared for exact checks
    \let\intnorm=\pgfmathresult
    \pgfmathparse{(\norm <= 4) && !(\m==0 && \n==0)}  % Symmetric radius filter, exclude origin
    \ifnum \pgfmathresult = 1
      \ifnum\intnorm=1  % Boundary (for r=1, intnorm=1)
        \fill[color=nearestgreen] (\x,\y) circle (2pt);
      \else
        \ifnum\intnorm<1  % Inner <1, but since inner is induced, skip originals inside (none <1 except origin)
        \else
          \fill[color=red] (\x,\y) circle (2pt);  % Outer: norm >1
          % Compute inverted blue point only for outer
          \pgfmathsetmacro{\xprime}{1 * \x / \intnorm}  % r^2=1 for r=1
          \pgfmathsetmacro{\yprime}{1 * \y / \intnorm}
          \pgfmathsetmacro{\distprime}{sqrt( (\xprime)*(\xprime) + (\yprime)*(\yprime) )}
          \pgfmathsetmacro{\radblue}{1 + ( (\distprime) - (1/sqrt(21)) ) / ( (1/sqrt(4)) - (1/sqrt(21)) ) }  % Adjusted scaling for r=1
          \fill[color=blue] (\xprime,\yprime) circle (0.75*\radblue pt);  % 25% smaller
        \fi
      \fi
    \fi
  }
}

% Example red path in S_2
% \draw[red, thick] (-2.5,{sqrt(3)/2}) -- (-1.5,{sqrt(3)/2}) -- (-2,{sqrt(3)});

% % Blue path (inverted)
% \draw[blue, thick] (-2.5/7, {sqrt(3)/14}) -- (-0.5, {sqrt(3)/6}) -- (-2/7, {sqrt(3)/7});
%
% % Gray dotted lines connecting vertices to their inverses
% \draw[gray, dotted, thick] (-2.5,{sqrt(3)/2}) -- (-2.5/7, {sqrt(3)/14});
% \draw[gray, dotted, thick] (-1.5,{sqrt(3)/2}) -- (-0.5, {sqrt(3)/6});
% \draw[gray, dotted, thick] (-2,{sqrt(3)}) -- (-2/7, {sqrt(3)/7});
%
% % Labels for red
% \node[below left, font=\small] at (-2.5,{sqrt(3)/2}) {$\vec{v}_1$};
% \node[below left, font=\small] at (-1.5,{sqrt(3)/2}) {$\vec{v}_2$};
% \node[below left, font=\small] at (-2,{sqrt(3)}) {$\vec{v}_3$};

% % Labels for blue
% \node[below left, font=\small] at (-0.45, 0.3) {$\iota_{1}(\vec{x})$};
% \node[below left, font=\small] at (-0.6, 0.6) {$\iota_{1}(\vec{y})$};
% \node[below left, font=\small] at (-0.25, 0.75) {$\iota_{1}(\vec{z})$};

\draw[dashed, color=nearestgreen] (1,0) -- (0.5,{sqrt(3)/2}) -- (-0.5,{sqrt(3)/2}) -- (-1,0) -- (-0.5,{-sqrt(3)/2}) -- (0.5,{-sqrt(3)/2}) -- cycle;  % Unit hexagon at primary directions

% Angular sector labels (adjusted for S_1 and S_4 to avoid axis intersection)
\foreach \t in {0,2,3,5} {
  \node[font=\Large] at ({\t*60 + 30}:3.9) {$S_\t$};
}
\node[font=\Large] at (85:3.6) {$S_1$}; % Shifted slightly right of 90°
\node[font=\Large] at (275:3.6) {$S_4$}; % Shifted slightly right of 270°

% Angle labels on axes (adjusted to avoid intersection)
\node[font=\Large, below right] at (3.5,0.45) {$\mathbb{R}$};
\node[font=\Large, above left] at (-0.16,3.6) {$\mathbb{I}$};

% Phase pair labels in quadrants
\node at (45:6) {Quadrant I: $(0, \pi/2)_\phi$};
\node at (135:6) {Quadrant II: $(\pi, \pi/2)_\phi$};
\node at (-135:6) {Quadrant III: $(\pi, 3\pi/2)_\phi$};
\node at (-45:6) {Quadrant IV: $(0, 3\pi/2)_\phi$};

% Axis phase labels
\node at (0:4.8) {East: $(0, 0)_\phi$};
\node at (90:4.25) {North: $(\pi/2, \pi/2)_\phi$};
\node at (180:4.8) {West: $(\pi, 0)_\phi$};
\node at (270:4.3) {South: $(\pi/2, 3\pi/2)_\phi$};
\end{tikzpicture}
}
\hspace*{-\dimexpr(\wd\figbox - \textwidth)/2 \relax}
\usebox\figbox

\caption{Zone and angular sector partitioning on a Tri-Quarter radial dual triangular lattice graph $\Lambda_1^4$ (with admissible inversion radius \(r = 1\) and truncation radius \(R=4\)). The boundary zone subgraph $\Lambda_{T,1}^4$ consists of six green vertices with uniform size that form a hexagon to encode an order-6 cyclic group, which align to precisely intersect the six angular sector boundary primary rays. The outer zone subgraph $\Lambda_{+,1}^4$ has red vertices with uniform size. The inner zone subgraph $\Lambda_{-,1}^4$ has blue vertices that decrease in (perceived) size as they approach the punctured origin (for illustration only). We see an example of Escher-inspired reflected twin paths: the red connected path in $\Lambda_{+,1}^4$ maps to the blue connected path in $\Lambda_{-,1}^4$ under circle inversion. For a dynamic version with randomly connected paths, see Appendix~\ref{appendix:radial-dual-lattice-graph} (with a screenshot in Figure~\ref{fig:zone-partition}).}
\label{fig:inverted-edge-original}
\vspace*{\fill}
\end{figure}